\chapter{A correct program can enforce the wrong rule}
\label{ch:problem}

% BEGIN SOURCE UNIT M01:00
\section{The decision before the decision}
\label{merge:M01:00}

\label{sec:problem-decision}
A distributor proposes a marketing offer involving an investment fund. The
offer includes an incentive. A compliance control must determine whether the
incentive falls within a restriction, whether an exception applies, and whether
the relevant evidence is sufficient to make that determination. A software team
can make these questions look simple: gather a few fields, evaluate a condition,
and return a result to the bank's Java system. The difficult decisions have
already occurred before the condition runs. Someone has decided what counts as
an incentive component, what the words in the circular cover, what evidence
establishes a promotional link, and which version of a referenced rule applies.

The distinction is consequential. Suppose a gift restriction covers both a
named investment product and a type of investment product. An implementation
that checks only the first route may work perfectly for a demonstration using
one named fund. Its unit tests may be comprehensive with respect to the code
that was written. Its developer may correctly explain every branch. The defect
is that a branch required by the adopted interpretation never reached the code.
Testing generated branches alone will not reveal a missing source obligation.

Now suppose that both promotional routes are represented, but an exception for
a fee discount is applied to an entire offer containing a fee waiver and a
separate voucher. The formal expression can be well typed, consistent and
executable. Its output can be reproduced by Python, Java and a solver. The
unresolved question is the object to which the exception attaches. A Boolean
variable named \texttt{discount} conceals that question unless its definition
identifies the assessed component and the evidence supporting the classification.

These are hypothetical failure scenarios derived from the retained 23EC46
exercise, not findings about a deployed bank system. They show how interpretation error can become a material operational risk. Regulatory action,
reputational damage and accountability for an erroneous control make the quality
of the decision process important. This book does not estimate the probability
or financial severity of an enforcement outcome. It specifies the evidence a
bank would need to understand, challenge and authorize its own interpretation.

The system we want to build therefore begins one step earlier than conventional
rule implementation. It asks what the bank has grounds to encode, which
alternative readings have been considered, what remains unknown, and who has
authority to resolve the uncertainty. Code generation follows that record. It
does not substitute for it.

The unit of work is a \emph{regulatory change assessment}. It includes source
material, a declared business scope, proposed interpretations, affected controls,
data requirements, implementation changes and review decisions. A single
circular can create several such assessments because different provisions may
concern different actors or activities. Conversely, one assessment may need
several circulars, Code provisions, FAQs and earlier decisions. Treating one
downloaded page as a self-contained legal specification would lose these
relationships.
% END SOURCE UNIT M01:00

% BEGIN SOURCE UNIT P-01-problem:00
\section{The bank workflow and its regulatory perimeter}
\label{merge:P-01-problem:00}

\label{sec:problem}

A private bank must turn regulatory requirements into client classifications,
product permissions, disclosures, monitoring, records and staff actions. Those
requirements rarely arrive as a complete algorithm. A circular may change one
part of an existing regime, depend on definitions in a code, distinguish product
providers from distributors, or explain an exception in an annex. Software must
nevertheless decide what to do for a particular client, product and event.

The sponsor's reported problem is the error-prone manual reading and updating of
bank systems. This proposal treats that account as the starting business need;
it does not assume a measured incident rate. The process below is a plausible
operating model to validate with the bank. A compliance officer reads a circular,
writes an interpretation, and discusses it with product and operations staff. A
business analyst turns the interpretation into requirements. Engineers implement
configuration or code, and testers construct examples from those requirements.
Each participant can perform competently while a condition disappears between
documents. When the circular changes, the bank must reconstruct the chain.

\begin{figure}[htbp]
\centering
\begin{tikzpicture}[node distance=5mm, every node/.style={font=\small},
  box/.style={draw=teal,rounded corners,fill=pale,align=center,text width=25mm,minimum height=15mm},
  arr/.style={-{Latex[length=2mm]},thick,draw=navy}]
\node[box] (a) {Circulars, annexes, codes};
\node[box,right=of a] (b) {Compliance interpretation};
\node[box,right=of b] (c) {Requirements and tickets};
\node[box,right=of c] (d) {System rules and workflow};
\node[box,right=of d] (e) {Client decisions and records};
\draw[arr] (a)--(b); \draw[arr] (b)--(c); \draw[arr] (c)--(d); \draw[arr] (d)--(e);
\end{tikzpicture}
\caption{The proposed discovery exercise will map the bank's actual process onto
these transitions. LegalMath preserves the source, interpretation and tests across
them, so a reviewer can follow a decision back to its reason.}
\label{fig:manual}
\end{figure}

The product opportunity is to make that chain inspectable and reusable. A
reviewer should be able to select a rule and see the passage that supports it,
the interpretation adopted, the facts it needs, the situations in which it changes
an outcome, and the version actually used. A subsequent amendment should identify
the affected controls and test cases without relying on the memory of the
original implementation team.
% END SOURCE UNIT P-01-problem:00

% BEGIN SOURCE UNIT P-01-problem:01
\section{The regulatory perimeter is part of the specification}
\label{merge:P-01-problem:01}


The supplied Hong Kong Securities and Futures Commission (SFC) product-authorization
page is a useful discovery source, but it
is not a complete inventory of a private bank's distribution obligations. The
SPI guidance examined here is a joint circular from the SFC and Hong Kong Monetary
Authority (HKMA) to intermediaries. The
tokenisation circular distinguishes product-provider responsibilities from
distributor responsibilities and imports other product and intermediary rules.
The marketing circular refers to the Securities and Futures Ordinance, the Code
of Conduct and advertising guidance. These dependencies determine the meaning of
the extracted clauses \citep{sfcindex,sfcspi,sfctoken2026,sfcmarketing}.

Discovery therefore starts with the bank's legal entities and activities. For a
specified business, the inventory should record whether the institution acts as
distributor, product provider, adviser or another regulated actor, and which
client and product classes are involved. Being a Hong Kong office of a US banking
group does not by itself settle any of these questions. The prototype will record
the selected Hong Kong perimeter and leave other group and jurisdictional
requirements as separately identified dependencies. A broader rollout requires
the bank's authoritative inventory.

Document classification also matters. The September 2025 joint circular announces
a concurrent thematic review of distribution practices and describes what the
review will examine. It should enter regulatory-change tracking; its description
of the review is not, by itself, a new numerical client-eligibility test
\citep{sfcreview2025}. A translator that treats every sentence as a binding
conditional will manufacture rules. Conversely, a translator that extracts only
sentences containing a particular modal verb can miss definitions, incorporated
requirements and material explanatory guidance.

LegalMath should preserve a passage's role: definition, requirement, permission,
exception, supervisory expectation, explanation or announcement. The original
wording and its authority remain alongside the adopted implementation. A bank may
choose an additional operational restriction, but that choice should be recorded
as bank policy with an owner and rationale. It must not be silently attributed to
the SFC.
% END SOURCE UNIT P-01-problem:01

% BEGIN SOURCE UNIT M01:01
\section{What the existing experiment actually established}
\label{merge:M01:01}

\label{sec:problem-evidence}
The present LegalMath workbench provides a useful starting point because it has
already crossed the engineering boundary from a formal rule to generated Java.
Its second-circular exercise concerns a deliberately narrow interpretation of
paragraph 10 of SFC circular 23EC46, dated 24 October 2023
\citep{sfcmarketing}. The example assesses one incentive component in a specified
distributor and SFC-authorised-fund profile. Whether the component is a gift,
whether it is a fee discount, and whether it has a relevant product link are
supplied classifications. The program does not establish those classifications
by reading an offer document.

The exercise preserved a packet of 22 named expected outcomes before constructing
the candidate rule. A separate truth-table implementation enumerated 729
combinations of true, false and unknown values across six inputs. The actual
generated Java class and the Python evaluator agreed on the declared results.
Four deliberately altered rules were compiled and detected by the named cases.
Those are substantial engineering observations: they show that the exercised
implementation computes the proposed rule in the checked domain, and that the
cases can expose several particular errors.

They do not provide independent evidence that the proposed rule is the correct
reading of the circular. The same implementing agent authored the interpretation
and the expected outcomes. Writing one before the other makes later alteration
detectable; it does not create a second legal judgment. The separate truth-table
oracle removes a direct dependency on the candidate's abstract syntax tree. It
still expresses the same chosen formula. If that choice omitted a legal
condition, all three implementations could agree on the omission.

The candidate also retained five unresolved review issues and two unresolved
source dependencies. Its attempt to enter review was rejected. It remained a
draft and no release was created. That refusal is part of the useful result.
The workbench did not mistake successful compilation for authority to deploy a
regulatory control. The monograph preserves this example as a record of an
unapproved candidate; resolving the cited dependencies here by assumption would
destroy the very distinction the example teaches.

The broader application acceptance contains 154 passing tests, but its drafting
provider returns supplied fixtures. It has not interpreted arbitrary circulars
in a measured live-model evaluation. The acceptance count is evidence about
source intake, data contracts, execution, review workflow and delivery under the
tested conditions. It is not an estimate of translation accuracy. A dashboard
that presents the count without the target of the tests would invite a stronger
conclusion than the evidence supports.

This distinction suggests the question the new product must answer. Given a
source packet and a proposed operational scope, can it expose a material
misinterpretation before that misinterpretation becomes an approved executable
control? Success requires more than generating valid syntax. It requires finding
missing alternatives, retrieving the context that changes a reading, and making
disagreements concrete enough for a competent reviewer to resolve.
% END SOURCE UNIT M01:01

% BEGIN SOURCE UNIT M01:02
\section{Five places where meaning can be lost}
\label{merge:M01:02}

\label{sec:problem-boundaries}
It is tempting to describe the workflow as one translation from English to
Java. That description hides several different kinds of error. The first is
source error: the reviewed text is incomplete, extracted incorrectly, superseded,
or wrongly treated as authoritative. An omitted footnote can remove a definition.
A copied paragraph can lose a heading that limits its scope. A source's
publication date can be mistaken for its effective date. None of these failures
requires a language model to reason incorrectly; the model may never receive
the material needed for the correct decision.

The second is interpretation error. The words are present, but their legal
function is misrepresented. A qualification becomes an unconditional rule. An
exception is attached to the wrong object. A statement about possible
classification becomes a finding of fact. A normative expectation is converted
into an operational block without recording the bank's choice. These are
decisions about meaning, not merely extraction or programming defects.

The third is specification error. The reviewer may have a defensible
interpretation, but the intermediate rule expresses something else. A conjunction
is replaced by a disjunction, a threshold loses equality, or an unknown condition
is interpreted as false. The specification can also fail through insufficient
expressiveness: a time-dependent obligation may be forced into a one-off Boolean
test. A software language that cannot represent a required distinction must
report the gap rather than approximate it silently.

The fourth is execution error. Java can evaluate a correct intermediate rule
incorrectly because of a compiler defect, numeric conversion, time-zone mismatch,
or inconsistent evidence handling. A data adapter may supply dollars where the
rule expects cents. A build may package an older runtime. An apparent match on a
few simple examples can miss a defect in unknown-value or conflict precedence.
These failures are often amenable to strong formal or exhaustive checks once
the specification and input domain are fixed.

The fifth is operational error. The correct Java result can be used for the
wrong client, product, source version or transaction time. Consent can change
between evaluation and order commitment. A system can treat the absence of one
prohibition as permission to proceed despite unresolved checks elsewhere. A
reviewed rule can remain active after its source is replaced. These are failures
of integration and change control, even when the individual rule evaluator is
correct.

Figure~\ref{fig:assurance-chain} gives each transition its own question. The
arrows do not imply that success flows automatically from one stage to the next. Every
arrow needs evidence appropriate to the claim it makes.

\begin{figure}[htbp]
\centering
\begin{tikzpicture}[node distance=7mm, every node/.style={font=\small},
box/.style={draw=navy,rounded corners,align=center,text width=125mm,minimum height=11mm},
arr/.style={-{Latex},draw=teal,thick}]
\node[box] (s) {Retained source: are these the complete, applicable words?};
\node[box,below=of s] (i) {Interpretation: what reading do the words and context support?};
\node[box,below=of i] (f) {Formal specification: does the rule express that reading?};
\node[box,below=of f] (j) {Java execution: does this binary preserve the rule?};
\node[box,below=of j] (a) {Bank action: is this result used with the right evidence and authority?};
\draw[arr] (s)--(i); \draw[arr] (i)--(f); \draw[arr] (f)--(j); \draw[arr] (j)--(a);
\end{tikzpicture}
\caption{A correct downstream step does not repair a mistaken upstream premise.
Each transition has a different assurance question.}
\label{fig:assurance-chain}
\end{figure}

A useful consequence follows immediately. Two checks count as meaningful
redundancy only to the extent that their opportunities for failure differ. A
second compiler can challenge code generation. It cannot detect a missing source
document if it receives the same incomplete specification. A second language
model may propose a different reading, but it cannot authenticate a source merely
by describing it confidently. Redundancy must be designed around the failure
being investigated.
% END SOURCE UNIT M01:02

% BEGIN SOURCE UNIT P-01-problem:05
\subsection{What makes the current work expensive and hard to verify}
\label{merge:P-01-problem:05}


The common failure is loss of context. Thresholds detach from definitions;
exceptions detach from their parent rules; duties detach from actors and events;
and historical decisions detach from the source version that justified them.
Document search helps locate text, but it does not preserve those relationships
through implementation. More fluent summaries can make an omitted condition
harder to notice.

Review then becomes repetitive. Compliance officers must infer what an engineer
understood, engineers must infer what a compliance memo leaves implicit, and
testers must construct enough cases to expose mismatches. If tests are written
only from the implementation ticket, a mistranslation shared by the ticket and
code can pass every test. After an amendment, teams may retest too little because
dependencies are hidden, or retest too much because they cannot show what is
unaffected.

The proposed product makes each material interpretation a maintained part of the
system. Its value should be measured in fewer material defects, more complete
explanations, and less repeated review effort. Those benefits are hypotheses for
the prototype to test. The evidence needed to test them includes independently
reviewed cases, disagreements that remain visible, and the time spent resolving
them; a high extraction score alone cannot establish the business case.
% END SOURCE UNIT P-01-problem:05

% BEGIN SOURCE UNIT M01:03
\section{The denominator problem in completeness}
\label{merge:M01:03}

\label{sec:problem-denominator}
Suppose a drafting agent extracts eight requirements from a circular and produces
eight rules. A coverage report then shows that every extracted requirement has a
rule, every rule has a test, and every test passes. The result sounds complete.
It is complete only relative to the agent's own extraction. If the circular
contains a ninth operative requirement, the report has no row on which that
omission can appear. The system has used its answer to define the denominator
of its completeness measure.

The source inventory must therefore originate independently of the proposed
rule set. At the document level it records paragraphs, headings, annexes,
footnotes, tables and incorporated references. At the normative level it records
potential duties, prohibitions, permissions, definitions, exceptions and
qualifications. These two inventories serve different purposes. A paragraph
inventory can establish that a passage was considered, but one paragraph may
contain several operative propositions. A rule inventory can establish that an
identified proposition has a disposition, but it cannot prove that the
propositions were all identified.

For the retained circular, the existing inventory contains 18 numbered
paragraphs and four footnotes. Only one paragraph has an executable control in
the second-circular example. Other passages concern context, other controls,
qualitative assessment or external dependencies. It would be wrong to describe
the circular as implemented because its selected gift rule passes 729 input
combinations. It would also be wrong to implement every paragraph mechanically:
an observation about market practices is not necessarily a new obligation, and a
qualitative duty may need a reviewed procedure rather than a fabricated numeric
threshold.

A complete disposition process gives every inventoried proposition a reasoned
destination. It may be linked to an executable rule, to a reviewed manual
procedure, to a justified exclusion from the declared assessment, or to an open
question. An open question is visible incompleteness. An exclusion requires a
scope rationale and reviewer authority. A manual procedure needs an owner and
evidence of performance; attaching the word ``manual'' to an inconvenient clause
does not implement it.

The same denominator problem occurs in search. If five candidate interpretations
agree, the correct conclusion is agreement among those five candidates under
their recorded assumptions. The system has not measured the interpretations it
failed to generate. Breadth constraints and targeted challenges improve the
opportunity to find missing readings, but there is no finite list of prompts
that proves unrestricted legal interpretation has been exhausted. The product
must report which dimensions were explored and which remain outside the search.

This makes completeness a collection of bounded claims. We can check that every
retained paragraph has a disposition. We can check that every cited reference has
been resolved or explicitly left open. We can check that each represented
exception has positive, negative and uncertain cases. We cannot infer from those
facts that no relevant circular exists elsewhere or that every legally
defensible reading has been found. A bank needs a process for setting and
maintaining its regulatory perimeter as well as a translator inside that perimeter.
% END SOURCE UNIT M01:03

% BEGIN SOURCE UNIT M01:04
\section{Why agreement is useful and insufficient}
\label{merge:M01:04}

\label{sec:problem-agreement}
An ensemble is a group of methods whose results are examined together. In a
prediction task it is common to combine answers by majority vote or averaging.
For regulatory interpretation, the useful unit of combination is often a
reasoned claim: one method identifies a scope restriction, another locates an
exception, and another generates a case that exposes a missing promotional route.
Their outputs need not be identical to contribute to the same assessment.

Agreement can provide evidence that a representation is stable across some
transformations. A controlled-English sentence and its formal rule may produce
the same decisions. A back-translation may preserve the formal structure. A
second implementation may agree with the first. Each observation rules out some
failure explanations. None rules out every explanation. In particular, methods
that share the same omitted premise can agree for the same wrong reason.

Disagreement is therefore not merely a nuisance to suppress. It may be the first
observable sign of a missing definition, a conflict in evidence, or a genuine
interpretive choice. If one candidate treats a product-type promotion as covered
and another does not, the difference points to a specific source phrase. If two
candidates disagree only when a classification is unknown, the problem may be
their evaluation semantics. If two sources have different effective intervals,
the disagreement may be temporal rather than substantive. The controller should
classify these possibilities before attempting a repair.

The converse is equally important: disagreement alone does not make every
candidate legally respectable. A generated rule that deletes an express exception
may be a useful negative control but a poor interpretation. The system should
retain it long enough to explain the rejection and test the detectors, while
distinguishing it from a reading that remains genuinely unresolved. A shortlist
filled with contrived alternatives can increase reviewer workload without
improving assurance.

The proposed ensemble consequently has two obligations. It must preserve
material dissent until there is a recorded reason to resolve it, and it must
identify the evidentiary standing of that dissent. An objection supported by an
applicable source is different from an unsupported model suggestion. A solver
counterexample proves a difference in formal consequences; it does not by itself
give one interpretation greater legal authority. A human reviewer can accept a
reading for a declared scope, but that decision must not be mislabelled as a
machine proof of English meaning.

Chapter~\ref{ch:ensemble} will specify how these distinctions become a finite
automatic process. The core idea is that a discrepancy creates work with a
defined purpose: retrieve missing authority, correct an encoding, investigate an
alternative or clarify a fact. Repetition without a new diagnostic objective is
not additional assurance. When the permitted actions have been exhausted, the
system must make the remaining issue easier for a reviewer to understand.
% END SOURCE UNIT M01:04

% BEGIN SOURCE UNIT M01:05
\section{The object being classified is part of the rule}
\label{merge:M01:05}

\label{sec:problem-objects}
The incentive example reveals a general weakness of compact decision tables:
they often define a field without defining the entity to which it belongs. A
field called \texttt{is\_discount} could describe a payment, a fee line, a
marketing offer, a product or a client relationship. These objects are not
interchangeable. If a fee waiver and voucher are separate incentive components,
a property of the waiver cannot automatically be transferred to the voucher.
The rule's subject identity is therefore part of its meaning.

The source-to-rule specification must state the unit of assessment before
introducing its predicates. In this example it is one proposed component within
a documented offer. The specification must also state how components are
identified and when a purported separation is legally meaningful. A data
designer cannot settle that second question solely by assigning two database
identifiers. Operational decomposition needs to reflect the reviewed legal
analysis; otherwise a bank could accidentally change the meaning of an exception
by changing how it stores an offer.

The same problem appears in private-bank suitability workflows. A statement
about a client's status does not necessarily apply to every account, every
product category or every transaction channel. Consent may be linked to an
identified service and period. A relationship manager's assessment may have a
subject, scope and evidence date that differ from those of a wealth record.
The intermediate representation must retain these bindings, not merely a
collection of client-level Boolean fields.

A useful specification sentence therefore contains both a predicate and its
subject: ``component C in offer O is classified as a fee discount under
definition D, on the evidence recorded in assessment A.'' That sentence is longer
than ``discount = true,'' but the extra information explains what the Boolean
means. The eventual runtime can carry compact identifiers so that ordinary
transactions remain efficient. The review record retains the definitions and
evidence those identifiers name.

This is also why a language model's ability to extract entities is insufficient.
It may correctly identify the words ``fee waiver'' and ``voucher'' while failing
to determine the legal relationship between them. Extraction can propose
candidates for classification. The system must separately record the reasoning
that accepts their relationship and the uncertainty that survives it. That
separation lets a later source clarification update a definition without
pretending that the underlying marketing document changed.
% END SOURCE UNIT M01:05

% BEGIN SOURCE UNIT M01:06
\section{Time changes the question being answered}
\label{merge:M01:06}

\label{sec:problem-time}
An apparently straightforward rule can yield different legitimate assessments
at different times. The regulation may change, the facts may change, or the
bank may learn something new about an earlier event. These are three distinct
sources of variation. A source publication date identifies when a document was
issued. An effective interval identifies when a reviewed rule is intended to
apply. An evidence recording time identifies when the bank had the relevant
information. Substituting one date for another changes the assessment.

Consider a gift classification entered on Tuesday for an offer made on Monday.
A historical review can ask what the bank should conclude now about Monday's
offer, using Tuesday's evidence. It can also ask what the system could conclude
on Monday from the information then recorded. Both questions are useful. They
require different knowledge cutoffs. Replaying the earlier decision with all
subsequently available facts would conceal the information gap that existed at
the time.

The current runtime makes assessment time and knowledge cutoff explicit through
\texttt{valid\_at} and \texttt{known\_at}. Known evidence has a validity interval
and a recording timestamp. An expired classification becomes unknown for an
assessment outside its interval. Evidence recorded after the historical cutoff
cannot be used in that historical assessment. A superseding record names the
record it corrects; simple arrival order does not determine which assertion is
authoritative.

These are engineering conventions selected for a precise replay question. They
do not determine the legal effect of every late discovery or correction. A
corrected factual record may establish that an obligation was performed on time,
while the bank's earlier inability to demonstrate that fact remains part of the
audit history. A regulatory amendment may apply prospectively, retrospectively or
with transition rules that need interpretation. The translator must represent
the relevant temporal question before the runtime can answer it.

The ensemble must therefore challenge time assumptions as well as the visible
substantive condition. One candidate can be correct under a historical source
version and wrong under the current one. Another can use a current definition
to interpret an earlier transaction without justification. A disagreement report
that omits the source revision and assessment times can make these distinct
questions look like a conflict about the same rule. Immutable context identifiers
are necessary to know what is actually being compared.
% END SOURCE UNIT M01:06

% BEGIN SOURCE UNIT M01:07
\section{From an answer to an accountable decision}
\label{merge:M01:07}

\label{sec:problem-accountability}
A useful product result contains more than a sentence announcing compliance or
non-compliance. It identifies the precise condition evaluated, the reviewed scope,
the facts and definitions used, the source version, the proposed interpretation,
the evidence of correct execution, and any unresolved dependency. These details
serve different readers. A compliance reviewer needs to see the interpretation.
A data owner needs to see the evidence mapping. An engineer needs reproducible
inputs and an artifact identity. A release owner needs to know which approvals
bind the exact binary being deployed.

The design should not force every reader to inspect every detail on every case.
The review interface can begin with the decision and the material reason, then
allow inspection of assumptions and counterexamples. The underlying record must
still support that inspection. A concise explanation is valuable when it is a
faithful view of a fuller record; it is dangerous when it replaces a missing
record with persuasive prose.

Accountability also requires distinguishing an interpretation from a bank policy.
A bank may decide to prohibit a class of offers while a regulatory question
remains unresolved. That can be a sensible operational choice under its own
authority. The system should record the policy owner, scope and reason for the
restriction. It should not attribute the restriction to an SFC sentence unless
the interpretation supports that attribution. Otherwise later reviewers cannot
tell whether a change requires new legal analysis or a policy decision by the
bank.

Likewise, a manual review route must have a defined effect. ``Refer to compliance''
is incomplete if no owner, service expectation, required evidence or transaction
state is specified. An unresolved assessment may hold an automated approval,
route the case for review or invoke an already approved conservative policy. The
choice belongs to the bank's operating design. The translator should never
invent a transaction permission merely because its retries have reached a limit.

The product's ambition is thus substantial but specific. It should improve the
quality of the questions reaching a human reviewer, preserve alternatives that
would otherwise disappear, and verify execution of the interpretation that is
eventually adopted. The next chapter follows two circulars through that route. The complete
SPI transaction makes execution concrete; the smaller marketing restriction
then exposes why independent source judgment, explicit scope and uncertainty
cannot be replaced by a large number of successful program evaluations.
% END SOURCE UNIT M01:07

